% 讲义主文件。编译：make all（先跑 symcheck 再出三版）| make light | make dark | make violet
\documentclass[UTF8,a4paper,11pt,fontset=none,oneside,openany]{ctexbook}
\setCJKmainfont{Noto Serif CJK SC}[BoldFont=Noto Serif CJK SC Bold,AutoFakeSlant=0.2]
\setCJKsansfont{Noto Sans CJK SC}[BoldFont=Noto Sans CJK SC Bold]
\setCJKmonofont{Noto Sans Mono CJK SC}[AutoFakeSlant=0.2]
\setmonofont{DejaVu Sans Mono}[Scale=0.85]
\usepackage[margin=24mm,headheight=16pt]{geometry}
\providecommand{\Edition}{light}   % light | dark | darkpurple，由 Makefile 传入
\usepackage[\Edition]{lacia}
\usepackage{lacia-symbols}   % 可选：符号旁注 + 符号表；不需要就删掉这两行和附录
% 符号登记表：旁注与附录“符号表”的唯一来源。
% \LaciaSymGroup{组键}{组标题}
% \LaciaNewSym{\宏}{键}{组}{宏的排版内容}{表中展示，空=同前}{旁注短释义，空=不加旁注}{表中长释义}
% \LaciaSymEntry{键}{组}{表中展示}{长释义}   只进符号表、不定义宏
% \LaciaSymWatch{裸写形式}{提示}             一符多义：symcheck 报告其裸写位置
% 检查：make symcheck
\LaciaSymGroup{search}{二分查找（第 1 章）}

\LaciaNewSym{\lbound}{lbound}{search}{\operatorname{lb}}{\operatorname{lb}(x)}{第一个 $\ge x$ 的下标}{升序数组中第一个不小于 $x$ 的元素下标；所有元素都小于 $x$ 时为 $n$}
\LaciaSymEntry{halfopen}{search}{[l,r)}{左闭右开区间：含 $l$，不含 $r$}

\linespread{1.2}\setlength{\parskip}{4pt}\raggedbottom
\setcounter{tocdepth}{1}
\hypersetup{pdftitle={二分查找讲义}}
\newcommand{\LaciaHeadLText}{二分查找}
\renewcommand{\LaciaHeadL}{\LaciaHeadLText}

\begin{document}
\frontmatter
\thispagestyle{empty}
\vspace*{38mm}\noindent{\fontsize{38}{50}\selectfont\bfseries\sffamily\color{LaciaDeep}二分查找\par}
\vspace{4mm}\noindent{\color{LaciaRule}\rule{\textwidth}{0.5pt}}
\newpage
\chapter*{如何使用这份讲义}
\addcontentsline{toc}{chapter}{如何使用这份讲义}
这份讲义只讲一件事：在一个有序数组里，怎样用比较次数尽量少的办法找到目标，以及为什么写对它比想象中难。全文只有一章，建议按顺序读完，再回头做章末的自测。

\textbf{全书脉络。}先说清二分查找成立的前提和它为什么快，再把朴素写法一步步改成不会出边界错误的写法，最后把``找一个值''推广为``找第一个满足条件的位置''。后面的内容都建立在这一推广之上。

\textbf{色块含义。}正文里有四种色块：
\begin{keypoint}值得记住的结论或判断方法。\end{keypoint}
\begin{pitfall}手算或编程时最容易错的地方。\end{pitfall}
\begin{errata}原始教材或课件中的错误及其修正。\end{errata}
\begin{supp}超出主线、但有助于理解的补充内容。\end{supp}

\textbf{符号。}专用记号在每节第一次出现时，右侧页边有简短旁注；完整含义见附录的符号表。

\tableofcontents
\mainmatter
\chapter{二分查找}
\begin{overview}
上一阶段我们用线性扫描在数组里找数，代价与长度成正比。数据一旦有序，这个代价可以降到对数级：每比较一次就能丢掉一半候选。本章要解决的问题是，如何把这个想法写成一段在任何边界输入下都正确的程序。你会先看到最常见的两种错误写法，再学到``半开区间 + 不变量''的写法，最后用同一套框架解决``第一个不小于目标的位置''这类变体。读完应当能不查资料写出无错的二分，并能说清每一行为什么这样写。
\end{overview}

\section{为什么有序就能更快}
设数组 $a$ 升序，长度为 $n$，目标为 $x$。取中点 $m$ 比较：若 $a[m]<x$，则 $x$ 只可能在右半边；若 $a[m]\ge x$，则只可能在左半边（含 $m$）。每次比较把候选区间减半，所以最多比较 $\lceil\log_2 n\rceil+1$ 次。

\begin{keypoint}
二分的本质不是``找中间那个数''，而是\textbf{维护一个不变量}：目标若存在，必在当前区间内。每一步都在保持这个不变量的前提下缩小区间。
\end{keypoint}

\section{把它写对}
下面的写法用左闭右开区间 $[l,r)$，不变量是：答案 $\nonote{\lbound}(x)$（第一个不小于 $x$ 的下标）总在 $[l,r]$ 内。

\begin{code}{lower\_bound.c}
int lower_bound(const int *a, int n, int x) {
    int l = 0, r = n;                 /* 答案范围 [l, r] */
    while (l < r) {
        int m = l + (r - l) / 2;      /* 不写 (l+r)/2，避免溢出 */
        if (a[m] < x) l = m + 1;      /* 答案在 m 右侧 */
        else          r = m;          /* m 可能就是答案 */
    }
    return l;                         /* l == r，即答案 */
}
\end{code}

\begin{pitfall}
最常见的三个错误：中点写成 $(l+r)/2$ 在大数组上溢出；循环条件写成 \tm{l <= r} 却用半开区间的更新方式，导致死循环；更新 \tm{r = m - 1} 把可能的答案丢掉。三者的共同原因都是没有先写出不变量。
\end{pitfall}

\begin{center}
\begin{tabularx}{.9\textwidth}{>{\raggedright\arraybackslash}p{3.2cm}Xl}
\toprule \lth{问题} & \lth{返回值} & \lth{调用方式}\\\midrule
第一个 $\ge x$ 的位置 & $l$ & \tm{lower\_bound(a,n,x)}\\
第一个 $>x$ 的位置 & $l$ & \tm{lower\_bound(a,n,x+1)}\\
元素 $x$ 的个数 & $\lbound(x+1)-\lbound(x)$ & 两次调用相减\\\bottomrule
\end{tabularx}
\end{center}

\begin{supp}
当``满足条件''是一个关于下标单调的判定（先全假、后全真），二分就可以用来求最小的真下标，与数组里存的是什么无关。这就是常说的``二分答案''。
\end{supp}

\section{自测}
\begin{enumerate}
\item 在不变量``答案在 $[l,r]$ 内''下，为什么循环条件是 $l<r$ 而不是 $l\le r$？
\item 若把 \tm{r = m} 改成 \tm{r = m - 1}，哪一类输入会出错？请构造一个最小反例。
\item 如何用一次二分判断 $x$ 是否在数组中？
\item 若数组改为降序，$\lbound(x)$ 应当怎样重新定义，代码要改哪一行？
\end{enumerate}

\appendix
\chapter{符号表}
正文每节第一次出现下列带旁注的符号时，右侧页边会有简短释义。这里给出全部记号的完整含义。
\LaciaSymbolTable
\end{document}
